\originalchapter{弱相互作用与电弱统一}\label{ch:weak}
\originalsection{带电流与低能极限}
弱带电流由 W 交换, 轻子顶点包含 $\gamma^\mu(1-\gamma^5)$. 当 $q^2\ll M_W^2$, 传播子可按 $1/M_W^2$ 展开, 得 Fermi 四费米作用, 系数满足 $G_F/\sqrt2=g^2/(8M_W^2)$. 这是低能有效理论, 当能量接近 W 质量时必须恢复传播子, 否则高能行为不可靠.

Muon 衰变 $\mu^-\to e^-+\bar\nu_e+\nu_\mu$ 的最低阶无质量电子近似给
\begin{equation}
 \Gamma_\mu=\frac{G_F^2m_\mu^5}{192\pi^3},\qquad
 \frac1\Gamma\frac{\dd\Gamma}{\dd x}=2x^2(3-2x),\quad x=\frac{2E_e}{m_\mu}.
\end{equation}
$m^5$ 可从四费米矩阵元与三体相空间的量纲先预见, $192\pi^3$ 则需实际自旋与相空间积分. 电子质量、辐射修正和极化关联会改变精确谱, 不应由最低阶数值冒称高精度测量结果.

中子衰变还包含核子矢量、轴矢量形状因子及 CKM 因子 $V_{ud}$. 能量释放仅约 $0.782\,\mathrm{MeV}$ 的电子动能, 与 muon 衰变尺度很不同. 允许近似的电子谱按式 \eqref{eq:beta-spectrum}, 还受质子库仑场修正, 不能用 $m_n^5$ 代入 muon 公式.

\originalsection{夸克混合与 CP 相位}
弱相互作用本征下型夸克是质量本征态的线性组合, $d'_i=V_{ij}d_j$. $V$ 为 CKM 幺正矩阵. 一般 $n\times n$ 幺正矩阵有 $n^2$ 个实参数: Hermitian 生成元有 $n$ 个实对角元和 $n(n-1)/2$ 个复非对角元. 实正交矩阵只有 $n(n-1)/2$ 个角.

夸克场可重定义相位, $2n-1$ 个相位不改变物理, 因而 CKM 剩 $n(n-1)/2$ 个角和 $(n-1)(n-2)/2$ 个不可去相位. 三代有三个角、一个 CP 相位; 两代没有此类物理 CP 相位. 幺正关系如 $V_{ud}V_{ub}^*+V_{cd}V_{cb}^*+V_{td}V_{tb}^*=0$ 在复平面形成三角形, 其面积与 CP 破坏相关.

\originalsection{中性流与混合角}
电弱规范结构是 SU(2)$_L\times$U(1)$_Y$, 取电荷关系 $Q=T_3+Y/2$. 规范场混合为
\begin{equation}
 A_\mu=\sin\theta_W W^3_\mu+\cos\theta_WB_\mu,
 \quad Z_\mu=\cos\theta_WW^3_\mu-\sin\theta_WB_\mu,
 \quad e=g\sin\theta_W=g'\cos\theta_W.
\end{equation}
Z 对费米子的矢量和轴矢量耦合为 $g_V=T_3-2Q\sin^2\theta_W$, $g_A=T_3$, 配以总体顶点因子 $g/(2\cos\theta_W)$. 光子只耦合电荷, 中微子虽然电中性仍有弱中性流. 中性流实验因此能把电弱作用同纯电磁作用区分.

\originalsection{Z 共振与不可见宽度}
在 Z 极点附近, 指定末态的峰截面近似
\begin{equation}
 \sigma_f^0=\frac{12\pi}{M_Z^2}\frac{\Gamma_e\Gamma_f}{\Gamma_Z^2}.
\end{equation}
扫能量给总宽度 $\Gamma_Z$, 可见轻子和强子道确定对应部分宽度. 总宽度减去可见部分便得 $\Gamma_{\rm inv}$, 再除理论计算的一代轻中微子宽度得到有效活跃代数. 它只数与 Z 耦合且质量低于 $M_Z/2$ 的开放态; 不排除很重的中微子或无 Z 耦合的 sterile 态. 这一间接测量不要求逐事件探测中微子.
